\documentclass[aspectratio=169,10pt]{beamer}
\usetheme{Madrid}
\usecolortheme{default}
\usepackage[utf8]{inputenc}
\usepackage[T1]{fontenc}
\usepackage{lmodern}
\usepackage{amsmath,amssymb}
\usepackage{enumitem}
\usepackage{xcolor}

\setbeamertemplate{navigation symbols}{}
\setbeamertemplate{footline}[frame number]
\setbeamertemplate{itemize items}[circle]
\setbeamerfont{frametitle}{size=\large,series=\bfseries}

\title{CIA Part 1 -- Fraud Risks}
\subtitle{100 Hard Questions with Detailed Explanations}
\author{Study Deck}
\date{\today}

\begin{document}
	
	\begin{frame}
		\titlepage
	\end{frame}
	
	\begin{frame}{Format \& Coverage}
		\small
		\begin{itemize}
			\item 100 hard questions on Fraud Risks.
			\item Each question is followed by a detailed explanation slide.
			\item Coverage: fraud definitions, fraud triangle, fraud diamond, red flags,
			occupational fraud, financial statement fraud, corruption, asset
			misappropriation, fraud risk assessment, fraud prevention/detection,
			fraud investigation, interviewing, evidence preservation, reporting,
			management override, collusion, IT fraud, and fraud-related standards.
		\end{itemize}
	\end{frame}
	
	% =========================================================
	% Q1
	% =========================================================
	\begin{frame}[t]{Q1}
		\small
		\textbf{Question:} Which of the following best defines fraud?
		\begin{enumerate}[label=\Alph*.,itemsep=1pt,topsep=2pt]
			\item An unintentional error in financial reporting.
			\item An intentional act involving deception to obtain an unjust advantage.
			\item A disagreement with management.
			\item A minor policy violation.
		\end{enumerate}
		\vspace{4pt}
		\textbf{Correct Answer: B}
	\end{frame}
	
	\begin{frame}[t]{Q1 -- Explanation}
		\small
		\textbf{Correct: B.} Fraud is an intentional act involving deception to obtain an unjust or illegal advantage. \\[4pt]
		\textbf{A wrong:} Errors are unintentional. \\[2pt]
		\textbf{C wrong:} Disagreement is not fraud. \\[2pt]
		\textbf{D wrong:} Minor policy violations may not be fraud.
	\end{frame}
	
	% =========================================================
	% Q2
	% =========================================================
	\begin{frame}[t]{Q2}
		\small
		\textbf{Question:} The fraud triangle consists of:
		\begin{enumerate}[label=\Alph*.,itemsep=1pt,topsep=2pt]
			\item Pressure, opportunity, rationalization.
			\item Risk, control, detection.
			\item Plan, do, check.
			\item Cash, inventory, receivables.
		\end{enumerate}
		\vspace{4pt}
		\textbf{Correct Answer: A}
	\end{frame}
	
	\begin{frame}[t]{Q2 -- Explanation}
		\small
		\textbf{Correct: A.} Cressey's fraud triangle: pressure (incentive), opportunity, and rationalization. \\[4pt]
		\textbf{B wrong:} That's the audit risk model. \\[2pt]
		\textbf{C wrong:} That's the Deming/PDCA cycle. \\[2pt]
		\textbf{D wrong:} Those are balance sheet accounts.
	\end{frame}
	
	% =========================================================
	% Q3
	% =========================================================
	\begin{frame}[t]{Q3}
		\small
		\textbf{Question:} The fraud diamond adds which element to the fraud triangle?
		\begin{enumerate}[label=\Alph*.,itemsep=1pt,topsep=2pt]
			\item Capability.
			\item Pressure.
			\item Opportunity.
			\item Rationalization.
		\end{enumerate}
		\vspace{4pt}
		\textbf{Correct Answer: A}
	\end{frame}
	
	\begin{frame}[t]{Q3 -- Explanation}
		\small
		\textbf{Correct: A.} The fraud diamond adds capability (the person's skills/position to commit fraud). \\[4pt]
		\textbf{B wrong:} Pressure is in the triangle. \\[2pt]
		\textbf{C wrong:} Opportunity is in the triangle. \\[2pt]
		\textbf{D wrong:} Rationalization is in the triangle.
	\end{frame}
	
	% =========================================================
	% Q4
	% =========================================================
	\begin{frame}[t]{Q4}
		\small
		\textbf{Question:} Which is NOT a category of occupational fraud?
		\begin{enumerate}[label=\Alph*.,itemsep=1pt,topsep=2pt]
			\item Asset misappropriation.
			\item Corruption.
			\item Financial statement fraud.
			\item Market manipulation.
		\end{enumerate}
		\vspace{4pt}
		\textbf{Correct Answer: D}
	\end{frame}
	
	\begin{frame}[t]{Q4 -- Explanation}
		\small
		\textbf{Correct: D.} ACFE's occupational fraud tree: asset misappropriation, corruption, financial statement fraud. Market manipulation is securities fraud, not occupational fraud. \\[4pt]
		\textbf{A wrong:} Asset misappropriation is a category. \\[2pt]
		\textbf{B wrong:} Corruption is a category. \\[2pt]
		\textbf{C wrong:} Financial statement fraud is a category.
	\end{frame}
	
	% =========================================================
	% Q5
	% =========================================================
	\begin{frame}[t]{Q5}
		\small
		\textbf{Question:} Which is the most common type of occupational fraud?
		\begin{enumerate}[label=\Alph*.,itemsep=1pt,topsep=2pt]
			\item Asset misappropriation.
			\item Corruption.
			\item Financial statement fraud.
			\item Market manipulation.
		\end{enumerate}
		\vspace{4pt}
		\textbf{Correct Answer: A}
	\end{frame}
	
	\begin{frame}[t]{Q5 -- Explanation}
		\small
		\textbf{Correct: A.} Asset misappropriation is the most common but least costly type of occupational fraud. \\[4pt]
		\textbf{B wrong:} Corruption is less common. \\[2pt]
		\textbf{C wrong:} Financial statement fraud is least common but most costly. \\[2pt]
		\textbf{D wrong:} Market manipulation is not occupational fraud.
	\end{frame}
	
	% =========================================================
	% Q6
	% =========================================================
	\begin{frame}[t]{Q6}
		\small
		\textbf{Question:} Which is the most costly type of occupational fraud?
		\begin{enumerate}[label=\Alph*.,itemsep=1pt,topsep=2pt]
			\item Asset misappropriation.
			\item Corruption.
			\item Financial statement fraud.
			\item Market manipulation.
		\end{enumerate}
		\vspace{4pt}
		\textbf{Correct Answer: C}
	\end{frame}
	
	\begin{frame}[t]{Q6 -- Explanation}
		\small
		\textbf{Correct: C.} Financial statement fraud is least common but most costly per case. \\[4pt]
		\textbf{A wrong:} Asset misappropriation is most common but least costly. \\[2pt]
		\textbf{B wrong:} Corruption is in the middle. \\[2pt]
		\textbf{D wrong:} Market manipulation is not occupational fraud.
	\end{frame}
	
	% =========================================================
	% Q7
	% =========================================================
	\begin{frame}[t]{Q7}
		\small
		\textbf{Question:} Which is a red flag for fraud?
		\begin{enumerate}[label=\Alph*.,itemsep=1pt,topsep=2pt]
			\item Strong segregation of duties.
			\item Lifestyle changes inconsistent with salary.
			\item Regular reconciliations.
			\item Timely reporting.
		\end{enumerate}
		\vspace{4pt}
		\textbf{Correct Answer: B}
	\end{frame}
	
	\begin{frame}[t]{Q7 -- Explanation}
		\small
		\textbf{Correct: B.} Unexplained lifestyle changes are a classic fraud red flag. \\[4pt]
		\textbf{A wrong:} Segregation of duties is a control. \\[2pt]
		\textbf{C wrong:} Reconciliations are controls. \\[2pt]
		\textbf{D wrong:} Timely reporting is a positive sign.
	\end{frame}
	
	% =========================================================
	% Q8
	% =========================================================
	\begin{frame}[t]{Q8}
		\small
		\textbf{Question:} Which is a red flag for financial statement fraud?
		\begin{enumerate}[label=\Alph*.,itemsep=1pt,topsep=2pt]
			\item Rapid revenue growth inconsistent with industry.
			\item Strong internal controls.
			\item Regular reconciliations.
			\item Timely reporting.
		\end{enumerate}
		\vspace{4pt}
		\textbf{Correct Answer: A}
	\end{frame}
	
	\begin{frame}[t]{Q8 -- Explanation}
		\small
		\textbf{Correct: A.} Revenue growth inconsistent with industry or cash flow is a red flag. \\[4pt]
		\textbf{B wrong:} Strong controls reduce risk. \\[2pt]
		\textbf{C wrong:} Reconciliations are controls. \\[2pt]
		\textbf{D wrong:} Timely reporting is a positive sign.
	\end{frame}
	
	% =========================================================
	% Q9
	% =========================================================
	\begin{frame}[t]{Q9}
		\small
		\textbf{Question:} Which is a red flag for asset misappropriation?
		\begin{enumerate}[label=\Alph*.,itemsep=1pt,topsep=2pt]
			\item Employee living beyond means.
			\item Strong segregation of duties.
			\item Regular reconciliations.
			\item Timely reporting.
		\end{enumerate}
		\vspace{4pt}
		\textbf{Correct Answer: A}
	\end{frame}
	
	\begin{frame}[t]{Q9 -- Explanation}
		\small
		\textbf{Correct: A.} Living beyond means is a classic red flag for asset misappropriation. \\[4pt]
		\textbf{B wrong:} Segregation of duties is a control. \\[2pt]
		\textbf{C wrong:} Reconciliations are controls. \\[2pt]
		\textbf{D wrong:} Timely reporting is a positive sign.
	\end{frame}
	
	% =========================================================
	% Q10
	% =========================================================
	\begin{frame}[t]{Q10}
		\small
		\textbf{Question:} Which is a red flag for corruption?
		\begin{enumerate}[label=\Alph*.,itemsep=1pt,topsep=2pt]
			\item Close relationship with a vendor.
			\item Strong segregation of duties.
			\item Regular reconciliations.
			\item Timely reporting.
		\end{enumerate}
		\vspace{4pt}
		\textbf{Correct Answer: A}
	\end{frame}
	
	\begin{frame}[t]{Q10 -- Explanation}
		\small
		\textbf{Correct: A.} Close relationships with vendors can indicate bribery or kickbacks. \\[4pt]
		\textbf{B wrong:} Segregation of duties is a control. \\[2pt]
		\textbf{C wrong:} Reconciliations are controls. \\[2pt]
		\textbf{D wrong:} Timely reporting is a positive sign.
	\end{frame}
	
	% =========================================================
	% Q11
	% =========================================================
	\begin{frame}[t]{Q11}
		\small
		\textbf{Question:} Which best describes "management override"?
		\begin{enumerate}[label=\Alph*.,itemsep=1pt,topsep=2pt]
			\item Management bypasses controls to commit fraud.
			\item Management follows all controls.
			\item Management documents controls.
			\item Management tests controls.
		\end{enumerate}
		\vspace{4pt}
		\textbf{Correct Answer: A}
	\end{frame}
	
	\begin{frame}[t]{Q11 -- Explanation}
		\small
		\textbf{Correct: A.} Management override is when management bypasses controls, often for fraud. It is a key limitation of internal control. \\[4pt]
		\textbf{B wrong:} Following controls is not override. \\[2pt]
		\textbf{C wrong:} Documentation is not override. \\[2pt]
		\textbf{D wrong:} Testing is not override.
	\end{frame}
	
	% =========================================================
	% Q12
	% =========================================================
	\begin{frame}[t]{Q12}
		\small
		\textbf{Question:} Which best describes "collusion"?
		\begin{enumerate}[label=\Alph*.,itemsep=1pt,topsep=2pt]
			\item Two or more people working together to commit fraud.
			\item One person acting alone.
			\item A control activity.
			\item A type of reconciliation.
		\end{enumerate}
		\vspace{4pt}
		\textbf{Correct Answer: A}
	\end{frame}
	
	\begin{frame}[t]{Q12 -- Explanation}
		\small
		\textbf{Correct: A.} Collusion is when two or more people work together to commit fraud, defeating segregation of duties. \\[4pt]
		\textbf{B wrong:} One person is not collusion. \\[2pt]
		\textbf{C wrong:} Collusion is not a control. \\[2pt]
		\textbf{D wrong:} Reconciliation is a control.
	\end{frame}
	
	% =========================================================
	% Q13
	% =========================================================
	\begin{frame}[t]{Q13}
		\small
		\textbf{Question:} Which best describes "fraud risk assessment"?
		\begin{enumerate}[label=\Alph*.,itemsep=1pt,topsep=2pt]
			\item A continuous process integrated into audit planning.
			\item A one-time exercise.
			\item Only the external auditor's responsibility.
			\item Only relevant for public companies.
		\end{enumerate}
		\vspace{4pt}
		\textbf{Correct Answer: A}
	\end{frame}
	
	\begin{frame}[t]{Q13 -- Explanation}
		\small
		\textbf{Correct: A.} Fraud risk assessment is continuous and integrated into audit planning. \\[4pt]
		\textbf{B wrong:} One-time assessments miss emerging risks. \\[2pt]
		\textbf{C wrong:} Internal audit also assesses fraud risk. \\[2pt]
		\textbf{D wrong:} All organizations face fraud risk.
	\end{frame}
	
	% =========================================================
	% Q14
	% =========================================================
	\begin{frame}[t]{Q14}
		\small
		\textbf{Question:} Which is the best way to prevent fraud?
		\begin{enumerate}[label=\Alph*.,itemsep=1pt,topsep=2pt]
			\item Strong internal controls and ethical culture.
			\item Ignoring red flags.
			\item Reducing segregation of duties.
			\item Eliminating reconciliations.
		\end{enumerate}
		\vspace{4pt}
		\textbf{Correct Answer: A}
	\end{frame}
	
	\begin{frame}[t]{Q14 -- Explanation}
		\small
		\textbf{Correct: A.} Fraud prevention relies on strong internal controls, ethical culture, and tone at the top. \\[4pt]
		\textbf{B wrong:} Ignoring red flags enables fraud. \\[2pt]
		\textbf{C wrong:} Reducing segregation increases risk. \\[2pt]
		\textbf{D wrong:} Eliminating reconciliations increases risk.
	\end{frame}
	
	% =========================================================
	% Q15
	% =========================================================
	\begin{frame}[t]{Q15}
		\small
		\textbf{Question:} Which is the best way to detect fraud?
		\begin{enumerate}[label=\Alph*.,itemsep=1pt,topsep=2pt]
			\item Tips and hotlines.
			\item Ignoring red flags.
			\item Reducing segregation of duties.
			\item Eliminating reconciliations.
		\end{enumerate}
		\vspace{4pt}
		\textbf{Correct Answer: A}
	\end{frame}
	
	\begin{frame}[t]{Q15 -- Explanation}
		\small
		\textbf{Correct: A.} Tips and hotlines are the most common way fraud is detected. \\[4pt]
		\textbf{B wrong:} Ignoring red flags enables fraud. \\[2pt]
		\textbf{C wrong:} Reducing segregation increases risk. \\[2pt]
		\textbf{D wrong:} Eliminating reconciliations increases risk.
	\end{frame}
	
	% =========================================================
	% Q16
	% =========================================================
	\begin{frame}[t]{Q16}
		\small
		\textbf{Question:} When fraud is suspected, the internal auditor should:
		\begin{enumerate}[label=\Alph*.,itemsep=1pt,topsep=2pt]
			\item Confront the suspect immediately.
			\item Notify appropriate authorities and recommend investigation.
			\item Ignore it until year-end.
			\item Destroy workpapers.
		\end{enumerate}
		\vspace{4pt}
		\textbf{Correct Answer: B}
	\end{frame}
	
	\begin{frame}[t]{Q16 -- Explanation}
		\small
		\textbf{Correct: B.} Standard 1210.A2: auditors must identify fraud indicators and report them to appropriate authorities; investigation is management/legal's job. \\[4pt]
		\textbf{A wrong:} Confronting the suspect can tip them off and destroy evidence. \\[2pt]
		\textbf{C wrong:} Delay allows fraud to continue. \\[2pt]
		\textbf{D wrong:} Destroying workpapers is unethical and illegal.
	\end{frame}
	
	% =========================================================
	% Q17
	% =========================================================
	\begin{frame}[t]{Q17}
		\small
		\textbf{Question:} Which best describes "fraud investigation"?
		\begin{enumerate}[label=\Alph*.,itemsep=1pt,topsep=2pt]
			\item A process to resolve allegations of fraud.
			\item A routine audit.
			\item A control activity.
			\item A reconciliation.
		\end{enumerate}
		\vspace{4pt}
		\textbf{Correct Answer: A}
	\end{frame}
	
	\begin{frame}[t]{Q17 -- Explanation}
		\small
		\textbf{Correct: A.} Fraud investigation resolves allegations of fraud through evidence gathering and interviews. \\[4pt]
		\textbf{B wrong:} Routine audit is separate. \\[2pt]
		\textbf{C wrong:} Investigation is not a control. \\[2pt]
		\textbf{D wrong:} Reconciliation is a control.
	\end{frame}
	
	% =========================================================
	% Q18
	% =========================================================
	\begin{frame}[t]{Q18}
		\small
		\textbf{Question:} Which is the best way to preserve fraud evidence?
		\begin{enumerate}[label=\Alph*.,itemsep=1pt,topsep=2pt]
			\item Document, secure, and maintain chain of custody.
			\item Delete suspicious files.
			\item Ignore evidence.
			\item Confront the suspect.
		\end{enumerate}
		\vspace{4pt}
		\textbf{Correct Answer: A}
	\end{frame}
	
	\begin{frame}[t]{Q18 -- Explanation}
		\small
		\textbf{Correct: A.} Evidence must be documented, secured, and chain of custody maintained for legal use. \\[4pt]
		\textbf{B wrong:} Deleting files destroys evidence. \\[2pt]
		\textbf{C wrong:} Ignoring evidence is unethical. \\[2pt]
		\textbf{D wrong:} Confronting can compromise the investigation.
	\end{frame}
	
	% =========================================================
	% Q19
	% =========================================================
	\begin{frame}[t]{Q19}
		\small
		\textbf{Question:} Which best describes "fraud interviewing"?
		\begin{enumerate}[label=\Alph*.,itemsep=1pt,topsep=2pt]
			\item A structured approach to gather information from interviewees.
			\item A confrontation.
			\item A control activity.
			\item A reconciliation.
		\end{enumerate}
		\vspace{4pt}
		\textbf{Correct Answer: A}
	\end{frame}
	
	\begin{frame}[t]{Q19 -- Explanation}
		\small
		\textbf{Correct: A.} Fraud interviewing gathers information through structured, non-accusatory questioning. \\[4pt]
		\textbf{B wrong:} Confrontation can be counterproductive. \\[2pt]
		\textbf{C wrong:} Interviewing is not a control. \\[2pt]
		\textbf{D wrong:} Reconciliation is a control.
	\end{frame}
	
	% =========================================================
	% Q20
	% =========================================================
	\begin{frame}[t]{Q20}
		\small
		\textbf{Question:} Which best describes "fraud reporting"?
		\begin{enumerate}[label=\Alph*.,itemsep=1pt,topsep=2pt]
			\item Communicating fraud findings to appropriate parties.
			\item Ignoring fraud.
			\item Confronting the suspect.
			\item Destroying evidence.
		\end{enumerate}
		\vspace{4pt}
		\textbf{Correct Answer: A}
	\end{frame}
	
	\begin{frame}[t]{Q20 -- Explanation}
		\small
		\textbf{Correct: A.} Fraud reporting communicates findings to management, the board, and possibly regulators/law enforcement. \\[4pt]
		\textbf{B wrong:} Ignoring is unethical. \\[2pt]
		\textbf{C wrong:} Confronting is not reporting. \\[2pt]
		\textbf{D wrong:} Destroying evidence is illegal.
	\end{frame}
	
	% =========================================================
	% Q21
	% =========================================================
	\begin{frame}[t]{Q21}
		\small
		\textbf{Question:} Which best describes "occupational fraud"?
		\begin{enumerate}[label=\Alph*.,itemsep=1pt,topsep=2pt]
			\item Fraud committed by employees against their employer.
			\item Fraud by external parties.
			\item Market manipulation.
			\item Tax evasion.
		\end{enumerate}
		\vspace{4pt}
		\textbf{Correct Answer: A}
	\end{frame}
	
	\begin{frame}[t]{Q21 -- Explanation}
		\small
		\textbf{Correct: A.} Occupational fraud is fraud committed by employees against their employer. \\[4pt]
		\textbf{B wrong:} External fraud is separate. \\[2pt]
		\textbf{C wrong:} Market manipulation is securities fraud. \\[2pt]
		\textbf{D wrong:} Tax evasion is a different crime.
	\end{frame}
	
	% =========================================================
	% Q22
	% =========================================================
	\begin{frame}[t]{Q22}
		\small
		\textbf{Question:} Which is a type of asset misappropriation?
		\begin{enumerate}[label=\Alph*.,itemsep=1pt,topsep=2pt]
			\item Skimming.
			\item Bribery.
			\item Financial statement fraud.
			\item Market manipulation.
		\end{enumerate}
		\vspace{4pt}
		\textbf{Correct Answer: A}
	\end{frame}
	
	\begin{frame}[t]{Q22 -- Explanation}
		\small
		\textbf{Correct: A.} Skimming (stealing cash before recording) is asset misappropriation. \\[4pt]
		\textbf{B wrong:} Bribery is corruption. \\[2pt]
		\textbf{C wrong:} Financial statement fraud is separate. \\[2pt]
		\textbf{D wrong:} Market manipulation is securities fraud.
	\end{frame}
	
	% =========================================================
	% Q23
	% =========================================================
	\begin{frame}[t]{Q23}
		\small
		\textbf{Question:} Which is a type of corruption?
		\begin{enumerate}[label=\Alph*.,itemsep=1pt,topsep=2pt]
			\item Bribery.
			\item Skimming.
			\item Financial statement fraud.
			\item Market manipulation.
		\end{enumerate}
		\vspace{4pt}
		\textbf{Correct Answer: A}
	\end{frame}
	
	\begin{frame}[t]{Q23 -- Explanation}
		\small
		\textbf{Correct: A.} Bribery is a type of corruption. \\[4pt]
		\textbf{B wrong:} Skimming is asset misappropriation. \\[2pt]
		\textbf{C wrong:} Financial statement fraud is separate. \\[2pt]
		\textbf{D wrong:} Market manipulation is securities fraud.
	\end{frame}
	
	% =========================================================
	% Q24
	% =========================================================
	\begin{frame}[t]{Q24}
		\small
		\textbf{Question:} Which is a type of financial statement fraud?
		\begin{enumerate}[label=\Alph*.,itemsep=1pt,topsep=2pt]
			\item Revenue recognition fraud.
			\item Skimming.
			\item Bribery.
			\item Market manipulation.
		\end{enumerate}
		\vspace{4pt}
		\textbf{Correct Answer: A}
	\end{frame}
	
	\begin{frame}[t]{Q24 -- Explanation}
		\small
		\textbf{Correct: A.} Revenue recognition fraud (e.g., premature or fictitious revenue) is financial statement fraud. \\[4pt]
		\textbf{B wrong:} Skimming is asset misappropriation. \\[2pt]
		\textbf{C wrong:} Bribery is corruption. \\[2pt]
		\textbf{D wrong:} Market manipulation is securities fraud.
	\end{frame}
	
	% =========================================================
	% Q25
	% =========================================================
	\begin{frame}[t]{Q25}
		\small
		\textbf{Question:} Which best describes "skimming"?
		\begin{enumerate}[label=\Alph*.,itemsep=1pt,topsep=2pt]
			\item Stealing cash before it is recorded.
			\item Stealing cash after it is recorded.
			\item Bribery.
			\item Financial statement fraud.
		\end{enumerate}
		\vspace{4pt}
		\textbf{Correct Answer: A}
	\end{frame}
	
	\begin{frame}[t]{Q25 -- Explanation}
		\small
		\textbf{Correct: A.} Skimming is stealing cash before it is recorded, making it harder to detect. \\[4pt]
		\textbf{B wrong:} Stealing after recording is larceny. \\[2pt]
		\textbf{C wrong:} Bribery is corruption. \\[2pt]
		\textbf{D wrong:} Financial statement fraud is separate.
	\end{frame}
	
	% =========================================================
	% Q26
	% =========================================================
	\begin{frame}[t]{Q26}
		\small
		\textbf{Question:} Which best describes "larceny"?
		\begin{enumerate}[label=\Alph*.,itemsep=1pt,topsep=2pt]
			\item Stealing cash after it is recorded.
			\item Stealing cash before it is recorded.
			\item Bribery.
			\item Financial statement fraud.
		\end{enumerate}
		\vspace{4pt}
		\textbf{Correct Answer: A}
	\end{frame}
	
	\begin{frame}[t]{Q26 -- Explanation}
		\small
		\textbf{Correct: A.} Larceny is stealing cash after it is recorded. \\[4pt]
		\textbf{B wrong:} Stealing before recording is skimming. \\[2pt]
		\textbf{C wrong:} Bribery is corruption. \\[2pt]
		\textbf{D wrong:} Financial statement fraud is separate.
	\end{frame}
	
	% =========================================================
	% Q27
	% =========================================================
	\begin{frame}[t]{Q27}
		\small
		\textbf{Question:} Which best describes "kiting"?
		\begin{enumerate}[label=\Alph*.,itemsep=1pt,topsep=2pt]
			\item Manipulating bank transfers to cover a cash shortage.
			\item Stealing cash before recording.
			\item Bribery.
			\item Financial statement fraud.
		\end{enumerate}
		\vspace{4pt}
		\textbf{Correct Answer: A}
	\end{frame}
	
	\begin{frame}[t]{Q27 -- Explanation}
		\small
		\textbf{Correct: A.} Kiting manipulates bank transfers (e.g., checks between accounts) to cover a cash shortage. \\[4pt]
		\textbf{B wrong:} Skimming is stealing before recording. \\[2pt]
		\textbf{C wrong:} Bribery is corruption. \\[2pt]
		\textbf{D wrong:} Financial statement fraud is separate.
	\end{frame}
	
	% =========================================================
	% Q28
	% =========================================================
	\begin{frame}[t]{Q28}
		\small
		\textbf{Question:} Which best describes "lapping"?
		\begin{enumerate}[label=\Alph*.,itemsep=1pt,topsep=2pt]
			\item Using subsequent receipts to cover a stolen receipt.
			\item Stealing cash before recording.
			\item Bribery.
			\item Financial statement fraud.
		\end{enumerate}
		\vspace{4pt}
		\textbf{Correct Answer: A}
	\end{frame}
	
	\begin{frame}[t]{Q28 -- Explanation}
		\small
		\textbf{Correct: A.} Lapping uses subsequent receipts to cover a stolen receipt, often in receivables. \\[4pt]
		\textbf{B wrong:} Skimming is stealing before recording. \\[2pt]
		\textbf{C wrong:} Bribery is corruption. \\[2pt]
		\textbf{D wrong:} Financial statement fraud is separate.
	\end{frame}
	
	% =========================================================
	% Q29
	% =========================================================
	\begin{frame}[t]{Q29}
		\small
		\textbf{Question:} Which best describes "ghost employee"?
		\begin{enumerate}[label=\Alph*.,itemsep=1pt,topsep=2pt]
			\item A fictitious employee on the payroll.
			\item A real employee.
			\item A vendor.
			\item A customer.
		\end{enumerate}
		\vspace{4pt}
		\textbf{Correct Answer: A}
	\end{frame}
	
	\begin{frame}[t]{Q29 -- Explanation}
		\small
		\textbf{Correct: A.} A ghost employee is a fictitious employee added to payroll to divert funds. \\[4pt]
		\textbf{B wrong:} Real employees are not ghosts. \\[2pt]
		\textbf{C wrong:} Vendors are separate. \\[2pt]
		\textbf{D wrong:} Customers are separate.
	\end{frame}
	
	% =========================================================
	% Q30
	% =========================================================
	\begin{frame}[t]{Q30}
		\small
		\textbf{Question:} Which best describes "kickback"?
		\begin{enumerate}[label=\Alph*.,itemsep=1pt,topsep=2pt]
			\item A payment to an employee for awarding a contract.
			\item A legitimate discount.
			\item A tax payment.
			\item A bonus.
		\end{enumerate}
		\vspace{4pt}
		\textbf{Correct Answer: A}
	\end{frame}
	
	\begin{frame}[t]{Q30 -- Explanation}
		\small
		\textbf{Correct: A.} A kickback is an illegal payment to an employee for awarding a contract or influencing a decision. \\[4pt]
		\textbf{B wrong:} Discounts are legitimate. \\[2pt]
		\textbf{C wrong:} Taxes are legal. \\[2pt]
		\textbf{D wrong:} Bonuses are legitimate.
	\end{frame}
	
	% =========================================================
	% Q31
	% =========================================================
	\begin{frame}[t]{Q31}
		\small
		\textbf{Question:} Which best describes "bribery"?
		\begin{enumerate}[label=\Alph*.,itemsep=1pt,topsep=2pt]
			\item Offering something of value to influence a decision.
			\item A legitimate discount.
			\item A tax payment.
			\item A bonus.
		\end{enumerate}
		\vspace{4pt}
		\textbf{Correct Answer: A}
	\end{frame}
	
	\begin{frame}[t]{Q31 -- Explanation}
		\small
		\textbf{Correct: A.} Bribery is offering something of value to influence a decision. \\[4pt]
		\textbf{B wrong:} Discounts are legitimate. \\[2pt]
		\textbf{C wrong:} Taxes are legal. \\[2pt]
		\textbf{D wrong:} Bonuses are legitimate.
	\end{frame}
	
	% =========================================================
	% Q32
	% =========================================================
	\begin{frame}[t]{Q32}
		\small
		\textbf{Question:} Which best describes "economic extortion"?
		\begin{enumerate}[label=\Alph*.,itemsep=1pt,topsep=2pt]
			\item Demanding payment to avoid harm.
			\item Offering a discount.
			\item Paying taxes.
			\item Receiving a bonus.
		\end{enumerate}
		\vspace{4pt}
		\textbf{Correct Answer: A}
	\end{frame}
	
	\begin{frame}[t]{Q32 -- Explanation}
		\small
		\textbf{Correct: A.} Economic extortion is demanding payment to avoid harm or unfavorable action. \\[4pt]
		\textbf{B wrong:} Discounts are legitimate. \\[2pt]
		\textbf{C wrong:} Taxes are legal. \\[2pt]
		\textbf{D wrong:} Bonuses are legitimate.
	\end{frame}
	
	% =========================================================
	% Q33
	% =========================================================
	\begin{frame}[t]{Q33}
		\small
		\textbf{Question:} Which best describes "conflict of interest"?
		\begin{enumerate}[label=\Alph*.,itemsep=1pt,topsep=2pt]
			\item A situation where personal interests could improperly influence judgment.
			\item A disagreement with management.
			\item A difference of opinion on audit scope.
			\item A budget dispute.
		\end{enumerate}
		\vspace{4pt}
		\textbf{Correct Answer: A}
	\end{frame}
	
	\begin{frame}[t]{Q33 -- Explanation}
		\small
		\textbf{Correct: A.} A conflict of interest exists when personal interests could improperly influence judgment. \\[4pt]
		\textbf{B wrong:} Disagreement is not automatically a conflict. \\[2pt]
		\textbf{C wrong:} Scope differences are professional. \\[2pt]
		\textbf{D wrong:} Budget disputes are administrative.
	\end{frame}
	
	% =========================================================
	% Q34
	% =========================================================
	\begin{frame}[t]{Q34}
		\small
		\textbf{Question:} Which best describes "financial statement fraud"?
		\begin{enumerate}[label=\Alph*.,itemsep=1pt,topsep=2pt]
			\item Intentional misstatement of financial statements.
			\item A clerical error.
			\item A disagreement with management.
			\item A budget dispute.
		\end{enumerate}
		\vspace{4pt}
		\textbf{Correct Answer: A}
	\end{frame}
	
	\begin{frame}[t]{Q34 -- Explanation}
		\small
		\textbf{Correct: A.} Financial statement fraud is intentional misstatement to deceive users. \\[4pt]
		\textbf{B wrong:} Clerical errors are unintentional. \\[2pt]
		\textbf{C wrong:} Disagreement is not fraud. \\[2pt]
		\textbf{D wrong:} Budget disputes are administrative.
	\end{frame}
	
	% =========================================================
	% Q35
	% =========================================================
	\begin{frame}[t]{Q35}
		\small
		\textbf{Question:} Which is a common financial statement fraud scheme?
		\begin{enumerate}[label=\Alph*.,itemsep=1pt,topsep=2pt]
			\item Premature revenue recognition.
			\item Skimming.
			\item Bribery.
			\item Ghost employees.
		\end{enumerate}
		\vspace{4pt}
		\textbf{Correct Answer: A}
	\end{frame}
	
	\begin{frame}[t]{Q35 -- Explanation}
		\small
		\textbf{Correct: A.} Premature revenue recognition is a common financial statement fraud scheme. \\[4pt]
		\textbf{B wrong:} Skimming is asset misappropriation. \\[2pt]
		\textbf{C wrong:} Bribery is corruption. \\[2pt]
		\textbf{D wrong:} Ghost employees are payroll fraud.
	\end{frame}
	
	% =========================================================
	% Q36
	% =========================================================
	\begin{frame}[t]{Q36}
		\small
		\textbf{Question:} Which is a common asset misappropriation scheme?
		\begin{enumerate}[label=\Alph*.,itemsep=1pt,topsep=2pt]
			\item Skimming.
			\item Premature revenue recognition.
			\item Bribery.
			\item Financial statement fraud.
		\end{enumerate}
		\vspace{4pt}
		\textbf{Correct Answer: A}
	\end{frame}
	
	\begin{frame}[t]{Q36 -- Explanation}
		\small
		\textbf{Correct: A.} Skimming is a common asset misappropriation scheme. \\[4pt]
		\textbf{B wrong:} Premature revenue is financial statement fraud. \\[2pt]
		\textbf{C wrong:} Bribery is corruption. \\[2pt]
		\textbf{D wrong:} Financial statement fraud is separate.
	\end{frame}
	
	% =========================================================
	% Q37
	% =========================================================
	\begin{frame}[t]{Q37}
		\small
		\textbf{Question:} Which is a common corruption scheme?
		\begin{enumerate}[label=\Alph*.,itemsep=1pt,topsep=2pt]
			\item Bribery.
			\item Skimming.
			\item Premature revenue recognition.
			\item Ghost employees.
		\end{enumerate}
		\vspace{4pt}
		\textbf{Correct Answer: A}
	\end{frame}
	
	\begin{frame}[t]{Q37 -- Explanation}
		\small
		\textbf{Correct: A.} Bribery is a common corruption scheme. \\[4pt]
		\textbf{B wrong:} Skimming is asset misappropriation. \\[2pt]
		\textbf{C wrong:} Premature revenue is financial statement fraud. \\[2pt]
		\textbf{D wrong:} Ghost employees are payroll fraud.
	\end{frame}
	
	% =========================================================
	% Q38
	% =========================================================
	\begin{frame}[t]{Q38}
		\small
		\textbf{Question:} Which best describes "fraud risk assessment"?
		\begin{enumerate}[label=\Alph*.,itemsep=1pt,topsep=2pt]
			\item Identifying fraud risks and evaluating controls.
			\item Ignoring fraud risks.
			\item Confronting suspects.
			\item Destroying evidence.
		\end{enumerate}
		\vspace{4pt}
		\textbf{Correct Answer: A}
	\end{frame}
	
	\begin{frame}[t]{Q38 -- Explanation}
		\small
		\textbf{Correct: A.} Fraud risk assessment identifies fraud risks and evaluates controls to mitigate them. \\[4pt]
		\textbf{B wrong:} Ignoring is unethical. \\[2pt]
		\textbf{C wrong:} Confronting is not assessment. \\[2pt]
		\textbf{D wrong:} Destroying evidence is illegal.
	\end{frame}
	
	% =========================================================
	% Q39
	% =========================================================
	\begin{frame}[t]{Q39}
		\small
		\textbf{Question:} Which best describes "fraud prevention"?
		\begin{enumerate}[label=\Alph*.,itemsep=1pt,topsep=2pt]
			\item Controls and culture to stop fraud before it occurs.
			\item Detecting fraud after it occurs.
			\item Confronting suspects.
			\item Destroying evidence.
		\end{enumerate}
		\vspace{4pt}
		\textbf{Correct Answer: A}
	\end{frame}
	
	\begin{frame}[t]{Q39 -- Explanation}
		\small
		\textbf{Correct: A.} Fraud prevention uses controls and culture to stop fraud before it occurs. \\[4pt]
		\textbf{B wrong:} Detection is after the fact. \\[2pt]
		\textbf{C wrong:} Confronting is not prevention. \\[2pt]
		\textbf{D wrong:} Destroying evidence is illegal.
	\end{frame}
	
	% =========================================================
	% Q40
	% =========================================================
	\begin{frame}[t]{Q40}
		\small
		\textbf{Question:} Which best describes "fraud detection"?
		\begin{enumerate}[label=\Alph*.,itemsep=1pt,topsep=2pt]
			\item Controls and procedures to identify fraud after it occurs.
			\item Controls to stop fraud before it occurs.
			\item Confronting suspects.
			\item Destroying evidence.
		\end{enumerate}
		\vspace{4pt}
		\textbf{Correct Answer: A}
	\end{frame}
	
	\begin{frame}[t]{Q40 -- Explanation}
		\small
		\textbf{Correct: A.} Fraud detection identifies fraud after it occurs (e.g., tips, reconciliations, analytics). \\[4pt]
		\textbf{B wrong:} Prevention is before the fact. \\[2pt]
		\textbf{C wrong:} Confronting is not detection. \\[2pt]
		\textbf{D wrong:} Destroying evidence is illegal.
	\end{frame}
	
	% =========================================================
	% Q41
	% =========================================================
	\begin{frame}[t]{Q41}
		\small
		\textbf{Question:} Which is the most common way fraud is detected?
		\begin{enumerate}[label=\Alph*.,itemsep=1pt,topsep=2pt]
			\item Tips and hotlines.
			\item External audit.
			\item Internal audit.
			\item Management review.
		\end{enumerate}
		\vspace{4pt}
		\textbf{Correct Answer: A}
	\end{frame}
	
	\begin{frame}[t]{Q41 -- Explanation}
		\small
		\textbf{Correct: A.} Tips and hotlines are the most common detection method. \\[4pt]
		\textbf{B wrong:} External audit detects some fraud. \\[2pt]
		\textbf{C wrong:} Internal audit detects some fraud. \\[2pt]
		\textbf{D wrong:} Management review detects some fraud.
	\end{frame}
	
	% =========================================================
	% Q42
	% =========================================================
	\begin{frame}[t]{Q42}
		\small
		\textbf{Question:} Which best describes "fraud risk factor"?
		\begin{enumerate}[label=\Alph*.,itemsep=1pt,topsep=2pt]
			\item A condition that increases the likelihood of fraud.
			\item A control activity.
			\item A reconciliation.
			\item A budget.
		\end{enumerate}
		\vspace{4pt}
		\textbf{Correct Answer: A}
	\end{frame}
	
	\begin{frame}[t]{Q42 -- Explanation}
		\small
		\textbf{Correct: A.} A fraud risk factor increases the likelihood of fraud (e.g., pressure, opportunity, rationalization). \\[4pt]
		\textbf{B wrong:} Control activities mitigate risk. \\[2pt]
		\textbf{C wrong:} Reconciliations are controls. \\[2pt]
		\textbf{D wrong:} Budgets are planning tools.
	\end{frame}
	
	% =========================================================
	% Q43
	% =========================================================
	\begin{frame}[t]{Q43}
		\small
		\textbf{Question:} Which best describes "fraud risk register"?
		\begin{enumerate}[label=\Alph*.,itemsep=1pt,topsep=2pt]
			\item A list of fraud risks and their controls.
			\item A list of employees.
			\item A list of vendors.
			\item A list of customers.
		\end{enumerate}
		\vspace{4pt}
		\textbf{Correct Answer: A}
	\end{frame}
	
	\begin{frame}[t]{Q43 -- Explanation}
		\small
		\textbf{Correct: A.} A fraud risk register documents fraud risks, their assessment, and controls. \\[4pt]
		\textbf{B wrong:} Employees are not a fraud register. \\[2pt]
		\textbf{C wrong:} Vendors are not a fraud register. \\[2pt]
		\textbf{D wrong:} Customers are not a fraud register.
	\end{frame}
	
	% =========================================================
	% Q44
	% =========================================================
	\begin{frame}[t]{Q44}
		\small
		\textbf{Question:} Which best describes "fraud risk assessment" in audit planning?
		\begin{enumerate}[label=\Alph*.,itemsep=1pt,topsep=2pt]
			\item Considering fraud risks when planning the engagement.
			\item Ignoring fraud risks.
			\item Confronting suspects.
			\item Destroying evidence.
		\end{enumerate}
		\vspace{4pt}
		\textbf{Correct Answer: A}
	\end{frame}
	
	\begin{frame}[t]{Q44 -- Explanation}
		\small
		\textbf{Correct: A.} Audit planning must consider fraud risks and design procedures to address them. \\[4pt]
		\textbf{B wrong:} Ignoring is unethical. \\[2pt]
		\textbf{C wrong:} Confronting is not planning. \\[2pt]
		\textbf{D wrong:} Destroying evidence is illegal.
	\end{frame}
	
	% =========================================================
	% Q45
	% =========================================================
	\begin{frame}[t]{Q45}
		\small
		\textbf{Question:} Which best describes "fraud brainstorming"?
		\begin{enumerate}[label=\Alph*.,itemsep=1pt,topsep=2pt]
			\item A team discussion of how fraud could occur.
			\item A control activity.
			\item A reconciliation.
			\item A budget.
		\end{enumerate}
		\vspace{4pt}
		\textbf{Correct Answer: A}
	\end{frame}
	
	\begin{frame}[t]{Q45 -- Explanation}
		\small
		\textbf{Correct: A.} Fraud brainstorming is a team discussion of how and where fraud could occur. \\[4pt]
		\textbf{B wrong:} It is not a control. \\[2pt]
		\textbf{C wrong:} It is not a reconciliation. \\[2pt]
		\textbf{D wrong:} It is not a budget.
	\end{frame}
	
	% =========================================================
	% Q46
	% =========================================================
	\begin{frame}[t]{Q46}
		\small
		\textbf{Question:} Which best describes "fraud risk factor -- pressure"?
		\begin{enumerate}[label=\Alph*.,itemsep=1pt,topsep=2pt]
			\item Financial or personal pressure on an individual.
			\item A control activity.
			\item A reconciliation.
			\item A budget.
		\end{enumerate}
		\vspace{4pt}
		\textbf{Correct Answer: A}
	\end{frame}
	
	\begin{frame}[t]{Q46 -- Explanation}
		\small
		\textbf{Correct: A.} Pressure is financial or personal stress that motivates fraud. \\[4pt]
		\textbf{B wrong:} It is not a control. \\[2pt]
		\textbf{C wrong:} It is not a reconciliation. \\[2pt]
		\textbf{D wrong:} It is not a budget.
	\end{frame}
	
	% =========================================================
	% Q47
	% =========================================================
	\begin{frame}[t]{Q47}
		\small
		\textbf{Question:} Which best describes "fraud risk factor -- opportunity"?
		\begin{enumerate}[label=\Alph*.,itemsep=1pt,topsep=2pt]
			\item Weak controls that allow fraud to occur.
			\item Strong controls.
			\item A reconciliation.
			\item A budget.
		\end{enumerate}
		\vspace{4pt}
		\textbf{Correct Answer: A}
	\end{frame}
	
	\begin{frame}[t]{Q47 -- Explanation}
		\small
		\textbf{Correct: A.} Opportunity is created by weak controls that allow fraud to occur. \\[4pt]
		\textbf{B wrong:} Strong controls reduce opportunity. \\[2pt]
		\textbf{C wrong:} Reconciliations are controls. \\[2pt]
		\textbf{D wrong:} Budgets are planning tools.
	\end{frame}
	
	% =========================================================
	% Q48
	% =========================================================
	\begin{frame}[t]{Q48}
		\small
		\textbf{Question:} Which best describes "fraud risk factor -- rationalization"?
		\begin{enumerate}[label=\Alph*.,itemsep=1pt,topsep=2pt]
			\item The fraudster justifies the act.
			\item Strong controls.
			\item A reconciliation.
			\item A budget.
		\end{enumerate}
		\vspace{4pt}
		\textbf{Correct Answer: A}
	\end{frame}
	
	\begin{frame}[t]{Q48 -- Explanation}
		\small
		\textbf{Correct: A.} Rationalization is how the fraudster justifies the act (e.g., "I was underpaid"). \\[4pt]
		\textbf{B wrong:} Strong controls reduce opportunity. \\[2pt]
		\textbf{C wrong:} Reconciliations are controls. \\[2pt]
		\textbf{D wrong:} Budgets are planning tools.
	\end{frame}
	
	% =========================================================
	% Q49
	% =========================================================
	\begin{frame}[t]{Q49}
		\small
		\textbf{Question:} Which best describes "fraud risk factor -- capability"?
		\begin{enumerate}[label=\Alph*.,itemsep=1pt,topsep=2pt]
			\item The person's skills and position to commit fraud.
			\item Strong controls.
			\item A reconciliation.
			\item A budget.
		\end{enumerate}
		\vspace{4pt}
		\textbf{Correct Answer: A}
	\end{frame}
	
	\begin{frame}[t]{Q49 -- Explanation}
		\small
		\textbf{Correct: A.} Capability is the person's skills, position, and ability to commit and conceal fraud. \\[4pt]
		\textbf{B wrong:} Strong controls reduce opportunity. \\[2pt]
		\textbf{C wrong:} Reconciliations are controls. \\[2pt]
		\textbf{D wrong:} Budgets are planning tools.
	\end{frame}
	
	% =========================================================
	% Q50
	% =========================================================
	\begin{frame}[t]{Q50}
		\small
		\textbf{Question:} Which best describes "fraud risk assessment" vs. "fraud investigation"?
		\begin{enumerate}[label=\Alph*.,itemsep=1pt,topsep=2pt]
			\item Assessment identifies risks; investigation resolves allegations.
			\item Assessment resolves; investigation identifies.
			\item They are the same.
			\item Neither matters.
		\end{enumerate}
		\vspace{4pt}
		\textbf{Correct Answer: A}
	\end{frame}
	
	\begin{frame}[t]{Q50 -- Explanation}
		\small
		\textbf{Correct: A.} Assessment identifies and evaluates fraud risks; investigation resolves allegations of fraud. \\[4pt]
		\textbf{B wrong:} Definitions are reversed. \\[2pt]
		\textbf{C wrong:} They are distinct. \\[2pt]
		\textbf{D wrong:} Both matter.
	\end{frame}
	
	% =========================================================
	% Q51
	% =========================================================
	\begin{frame}[t]{Q51}
		\small
		\textbf{Question:} Which best describes "fraud prevention" vs. "fraud detection"?
		\begin{enumerate}[label=\Alph*.,itemsep=1pt,topsep=2pt]
			\item Prevention stops fraud; detection finds fraud.
			\item Prevention finds; detection stops.
			\item They are the same.
			\item Neither matters.
		\end{enumerate}
		\vspace{4pt}
		\textbf{Correct Answer: A}
	\end{frame}
	
	\begin{frame}[t]{Q51 -- Explanation}
		\small
		\textbf{Correct: A.} Prevention stops fraud before it occurs; detection finds it after. \\[4pt]
		\textbf{B wrong:} Definitions are reversed. \\[2pt]
		\textbf{C wrong:} They are distinct. \\[2pt]
		\textbf{D wrong:} Both matter.
	\end{frame}
	
	% =========================================================
	% Q52
	% =========================================================
	\begin{frame}[t]{Q52}
		\small
		\textbf{Question:} Which best describes "fraud risk assessment" vs. "fraud risk register"?
		\begin{enumerate}[label=\Alph*.,itemsep=1pt,topsep=2pt]
			\item Assessment is the process; register documents the results.
			\item Register is the process; assessment documents.
			\item They are the same.
			\item Neither matters.
		\end{enumerate}
		\vspace{4pt}
		\textbf{Correct Answer: A}
	\end{frame}
	
	\begin{frame}[t]{Q52 -- Explanation}
		\small
		\textbf{Correct: A.} Fraud risk assessment is the process; the fraud risk register documents the results. \\[4pt]
		\textbf{B wrong:} Definitions are reversed. \\[2pt]
		\textbf{C wrong:} They are distinct. \\[2pt]
		\textbf{D wrong:} Both matter.
	\end{frame}
	
	% =========================================================
	% Q53
	% =========================================================
	\begin{frame}[t]{Q53}
		\small
		\textbf{Question:} Which best describes "fraud risk factor" vs. "fraud risk"?
		\begin{enumerate}[label=\Alph*.,itemsep=1pt,topsep=2pt]
			\item Factor increases likelihood; risk is the potential event.
			\item Risk increases likelihood; factor is the event.
			\item They are the same.
			\item Neither matters.
		\end{enumerate}
		\vspace{4pt}
		\textbf{Correct Answer: A}
	\end{frame}
	
	\begin{frame}[t]{Q53 -- Explanation}
		\small
		\textbf{Correct: A.} A fraud risk factor increases the likelihood of fraud; fraud risk is the potential event itself. \\[4pt]
		\textbf{B wrong:} Definitions are reversed. \\[2pt]
		\textbf{C wrong:} They are distinct. \\[2pt]
		\textbf{D wrong:} Both matter.
	\end{frame}
	
	% =========================================================
	% Q54
	% =========================================================
	\begin{frame}[t]{Q54}
		\small
		\textbf{Question:} Which best describes "fraud risk assessment" vs. "risk assessment"?
		\begin{enumerate}[label=\Alph*.,itemsep=1pt,topsep=2pt]
			\item Fraud risk assessment focuses on intentional acts; risk assessment is broader.
			\item They are the same.
			\item Risk assessment focuses on fraud only.
			\item Neither matters.
		\end{enumerate}
		\vspace{4pt}
		\textbf{Correct Answer: A}
	\end{frame}
	
	\begin{frame}[t]{Q54 -- Explanation}
		\small
		\textbf{Correct: A.} Fraud risk assessment focuses on intentional acts; general risk assessment covers all risks. \\[4pt]
		\textbf{B wrong:} They differ in scope. \\[2pt]
		\textbf{C wrong:} General risk assessment is broader. \\[2pt]
		\textbf{D wrong:} Both matter.
	\end{frame}
	
	% =========================================================
	% Q55
	% =========================================================
	\begin{frame}[t]{Q55}
		\small
		\textbf{Question:} Which best describes "fraud risk assessment" vs. "control risk assessment"?
		\begin{enumerate}[label=\Alph*.,itemsep=1pt,topsep=2pt]
			\item Fraud risk assessment focuses on intentional acts; control risk assessment focuses on control failures.
			\item They are the same.
			\item Control risk assessment focuses on fraud only.
			\item Neither matters.
		\end{enumerate}
		\vspace{4pt}
		\textbf{Correct Answer: A}
	\end{frame}
	
	\begin{frame}[t]{Q55 -- Explanation}
		\small
		\textbf{Correct: A.} Fraud risk assessment focuses on intentional acts; control risk assessment focuses on control failures. \\[4pt]
		\textbf{B wrong:} They differ. \\[2pt]
		\textbf{C wrong:} Control risk assessment is broader. \\[2pt]
		\textbf{D wrong:} Both matter.
	\end{frame}
	
	% =========================================================
	% Q56
	% =========================================================
	\begin{frame}[t]{Q56}
		\small
		\textbf{Question:} Which best describes "fraud risk assessment" vs. "inherent risk assessment"?
		\begin{enumerate}[label=\Alph*.,itemsep=1pt,topsep=2pt]
			\item Fraud risk assessment focuses on intentional acts; inherent risk is before controls.
			\item They are the same.
			\item Inherent risk focuses on fraud only.
			\item Neither matters.
		\end{enumerate}
		\vspace{4pt}
		\textbf{Correct Answer: A}
	\end{frame}
	
	\begin{frame}[t]{Q56 -- Explanation}
		\small
		\textbf{Correct: A.} Fraud risk assessment focuses on intentional acts; inherent risk is risk before controls. \\[4pt]
		\textbf{B wrong:} They differ. \\[2pt]
		\textbf{C wrong:} Inherent risk is broader. \\[2pt]
		\textbf{D wrong:} Both matter.
	\end{frame}
	
	% =========================================================
	% Q57
	% =========================================================
	\begin{frame}[t]{Q57}
		\small
		\textbf{Question:} Which best describes "fraud risk assessment" vs. "residual risk assessment"?
		\begin{enumerate}[label=\Alph*.,itemsep=1pt,topsep=2pt]
			\item Fraud risk assessment focuses on intentional acts; residual risk is after controls.
			\item They are the same.
			\item Residual risk focuses on fraud only.
			\item Neither matters.
		\end{enumerate}
		\vspace{4pt}
		\textbf{Correct Answer: A}
	\end{frame}
	
	\begin{frame}[t]{Q57 -- Explanation}
		\small
		\textbf{Correct: A.} Fraud risk assessment focuses on intentional acts; residual risk is risk after controls. \\[4pt]
		\textbf{B wrong:} They differ. \\[2pt]
		\textbf{C wrong:} Residual risk is broader. \\[2pt]
		\textbf{D wrong:} Both matter.
	\end{frame}
	
	% =========================================================
	% Q58
	% =========================================================
	\begin{frame}[t]{Q58}
		\small
		\textbf{Question:} Which best describes "fraud risk assessment" vs. "fraud risk response"?
		\begin{enumerate}[label=\Alph*.,itemsep=1pt,topsep=2pt]
			\item Assessment identifies risks; response treats them.
			\item Response identifies; assessment treats.
			\item They are the same.
			\item Neither matters.
		\end{enumerate}
		\vspace{4pt}
		\textbf{Correct Answer: A}
	\end{frame}
	
	\begin{frame}[t]{Q58 -- Explanation}
		\small
		\textbf{Correct: A.} Assessment identifies and evaluates fraud risks; response treats them (avoid, reduce, share, accept). \\[4pt]
		\textbf{B wrong:} Definitions are reversed. \\[2pt]
		\textbf{C wrong:} They are distinct. \\[2pt]
		\textbf{D wrong:} Both matter.
	\end{frame}
	
	% =========================================================
	% Q59
	% =========================================================
	\begin{frame}[t]{Q59}
		\small
		\textbf{Question:} Which best describes "fraud risk assessment" vs. "fraud risk monitoring"?
		\begin{enumerate}[label=\Alph*.,itemsep=1pt,topsep=2pt]
			\item Assessment identifies risks; monitoring tracks them over time.
			\item Monitoring identifies; assessment tracks.
			\item They are the same.
			\item Neither matters.
		\end{enumerate}
		\vspace{4pt}
		\textbf{Correct Answer: A}
	\end{frame}
	
	\begin{frame}[t]{Q59 -- Explanation}
		\small
		\textbf{Correct: A.} Assessment identifies fraud risks; monitoring tracks them and the effectiveness of controls over time. \\[4pt]
		\textbf{B wrong:} Definitions are reversed. \\[2pt]
		\textbf{C wrong:} They are distinct. \\[2pt]
		\textbf{D wrong:} Both matter.
	\end{frame}
	
	% =========================================================
	% Q60
	% =========================================================
	\begin{frame}[t]{Q60}
		\small
		\textbf{Question:} Which best describes "fraud risk assessment" vs. "fraud risk communication"?
		\begin{enumerate}[label=\Alph*.,itemsep=1pt,topsep=2pt]
			\item Assessment identifies risks; communication reports them.
			\item Communication identifies; assessment reports.
			\item They are the same.
			\item Neither matters.
		\end{enumerate}
		\vspace{4pt}
		\textbf{Correct Answer: A}
	\end{frame}
	
	\begin{frame}[t]{Q60 -- Explanation}
		\small
		\textbf{Correct: A.} Assessment identifies fraud risks; communication reports them to stakeholders. \\[4pt]
		\textbf{B wrong:} Definitions are reversed. \\[2pt]
		\textbf{C wrong:} They are distinct. \\[2pt]
		\textbf{D wrong:} Both matter.
	\end{frame}
	
	% =========================================================
	% Q61
	% =========================================================
	\begin{frame}[t]{Q61}
		\small
		\textbf{Question:} Which best describes "fraud risk assessment" vs. "fraud risk governance"?
		\begin{enumerate}[label=\Alph*.,itemsep=1pt,topsep=2pt]
			\item Assessment identifies risks; governance oversees them.
			\item Governance identifies; assessment oversees.
			\item They are the same.
			\item Neither matters.
		\end{enumerate}
		\vspace{4pt}
		\textbf{Correct Answer: A}
	\end{frame}
	
	\begin{frame}[t]{Q61 -- Explanation}
		\small
		\textbf{Correct: A.} Assessment identifies fraud risks; governance oversees their management. \\[4pt]
		\textbf{B wrong:} Definitions are reversed. \\[2pt]
		\textbf{C wrong:} They are distinct. \\[2pt]
		\textbf{D wrong:} Both matter.
	\end{frame}
	
	% =========================================================
	% Q62
	% =========================================================
	\begin{frame}[t]{Q62}
		\small
		\textbf{Question:} Which best describes "fraud risk assessment" vs. "fraud risk culture"?
		\begin{enumerate}[label=\Alph*.,itemsep=1pt,topsep=2pt]
			\item Assessment identifies risks; culture shapes behavior.
			\item Culture identifies; assessment shapes.
			\item They are the same.
			\item Neither matters.
		\end{enumerate}
		\vspace{4pt}
		\textbf{Correct Answer: A}
	\end{frame}
	
	\begin{frame}[t]{Q62 -- Explanation}
		\small
		\textbf{Correct: A.} Assessment identifies fraud risks; culture shapes how people behave toward risk. \\[4pt]
		\textbf{B wrong:} Definitions are reversed. \\[2pt]
		\textbf{C wrong:} They are distinct. \\[2pt]
		\textbf{D wrong:} Both matter.
	\end{frame}
	
	% =========================================================
	% Q63
	% =========================================================
	\begin{frame}[t]{Q63}
		\small
		\textbf{Question:} Which best describes "fraud risk assessment" vs. "fraud risk appetite"?
		\begin{enumerate}[label=\Alph*.,itemsep=1pt,topsep=2pt]
			\item Assessment identifies risks; appetite is willingness to accept them.
			\item Appetite identifies; assessment accepts.
			\item They are the same.
			\item Neither matters.
		\end{enumerate}
		\vspace{4pt}
		\textbf{Correct Answer: A}
	\end{frame}
	
	\begin{frame}[t]{Q63 -- Explanation}
		\small
		\textbf{Correct: A.} Assessment identifies fraud risks; appetite is the organization's willingness to accept fraud risk. \\[4pt]
		\textbf{B wrong:} Definitions are reversed. \\[2pt]
		\textbf{C wrong:} They are distinct. \\[2pt]
		\textbf{D wrong:} Both matter.
	\end{frame}
	
	% =========================================================
	% Q64
	% =========================================================
	\begin{frame}[t]{Q64}
		\small
		\textbf{Question:} Which best describes "fraud risk assessment" vs. "fraud risk tolerance"?
		\begin{enumerate}[label=\Alph*.,itemsep=1pt,topsep=2pt]
			\item Assessment identifies risks; tolerance is acceptable variation.
			\item Tolerance identifies; assessment accepts.
			\item They are the same.
			\item Neither matters.
		\end{enumerate}
		\vspace{4pt}
		\textbf{Correct Answer: A}
	\end{frame}
	
	\begin{frame}[t]{Q64 -- Explanation}
		\small
		\textbf{Correct: A.} Assessment identifies fraud risks; tolerance is acceptable variation around appetite. \\[4pt]
		\textbf{B wrong:} Definitions are reversed. \\[2pt]
		\textbf{C wrong:} They are distinct. \\[2pt]
		\textbf{D wrong:} Both matter.
	\end{frame}
	
	% =========================================================
	% Q65
	% =========================================================
	\begin{frame}[t]{Q65}
		\small
		\textbf{Question:} Which best describes "fraud risk assessment" vs. "fraud risk response"?
		\begin{enumerate}[label=\Alph*.,itemsep=1pt,topsep=2pt]
			\item Assessment identifies risks; response treats them.
			\item Response identifies; assessment treats.
			\item They are the same.
			\item Neither matters.
		\end{enumerate}
		\vspace{4pt}
		\textbf{Correct Answer: A}
	\end{frame}
	
	\begin{frame}[t]{Q65 -- Explanation}
		\small
		\textbf{Correct: A.} Assessment identifies fraud risks; response treats them. \\[4pt]
		\textbf{B wrong:} Definitions are reversed. \\[2pt]
		\textbf{C wrong:} They are distinct. \\[2pt]
		\textbf{D wrong:} Both matter.
	\end{frame}
	
	% =========================================================
	% Q66
	% =========================================================
	\begin{frame}[t]{Q66}
		\small
		\textbf{Question:} Which best describes "fraud risk assessment" vs. "fraud risk monitoring"?
		\begin{enumerate}[label=\Alph*.,itemsep=1pt,topsep=2pt]
			\item Assessment identifies risks; monitoring tracks them.
			\item Monitoring identifies; assessment tracks.
			\item They are the same.
			\item Neither matters.
		\end{enumerate}
		\vspace{4pt}
		\textbf{Correct Answer: A}
	\end{frame}
	
	\begin{frame}[t]{Q66 -- Explanation}
		\small
		\textbf{Correct: A.} Assessment identifies fraud risks; monitoring tracks them over time. \\[4pt]
		\textbf{B wrong:} Definitions are reversed. \\[2pt]
		\textbf{C wrong:} They are distinct. \\[2pt]
		\textbf{D wrong:} Both matter.
	\end{frame}
	
	% =========================================================
	% Q67
	% =========================================================
	\begin{frame}[t]{Q67}
		\small
		\textbf{Question:} Which best describes "fraud risk assessment" vs. "fraud risk communication"?
		\begin{enumerate}[label=\Alph*.,itemsep=1pt,topsep=2pt]
			\item Assessment identifies risks; communication reports them.
			\item Communication identifies; assessment reports.
			\item They are the same.
			\item Neither matters.
		\end{enumerate}
		\vspace{4pt}
		\textbf{Correct Answer: A}
	\end{frame}
	
	\begin{frame}[t]{Q67 -- Explanation}
		\small
		\textbf{Correct: A.} Assessment identifies fraud risks; communication reports them. \\[4pt]
		\textbf{B wrong:} Definitions are reversed. \\[2pt]
		\textbf{C wrong:} They are distinct. \\[2pt]
		\textbf{D wrong:} Both matter.
	\end{frame}
	
	% =========================================================
	% Q68
	% =========================================================
	\begin{frame}[t]{Q68}
		\small
		\textbf{Question:} Which best describes "fraud risk assessment" vs. "fraud risk governance"?
		\begin{enumerate}[label=\Alph*.,itemsep=1pt,topsep=2pt]
			\item Assessment identifies risks; governance oversees them.
			\item Governance identifies; assessment oversees.
			\item They are the same.
			\item Neither matters.
		\end{enumerate}
		\vspace{4pt}
		\textbf{Correct Answer: A}
	\end{frame}
	
	\begin{frame}[t]{Q68 -- Explanation}
		\small
		\textbf{Correct: A.} Assessment identifies fraud risks; governance oversees their management. \\[4pt]
		\textbf{B wrong:} Definitions are reversed. \\[2pt]
		\textbf{C wrong:} They are distinct. \\[2pt]
		\textbf{D wrong:} Both matter.
	\end{frame}
	
	% =========================================================
	% Q69
	% =========================================================
	\begin{frame}[t]{Q69}
		\small
		\textbf{Question:} Which best describes "fraud risk assessment" vs. "fraud risk culture"?
		\begin{enumerate}[label=\Alph*.,itemsep=1pt,topsep=2pt]
			\item Assessment identifies risks; culture shapes behavior.
			\item Culture identifies; assessment shapes.
			\item They are the same.
			\item Neither matters.
		\end{enumerate}
		\vspace{4pt}
		\textbf{Correct Answer: A}
	\end{frame}
	
	\begin{frame}[t]{Q69 -- Explanation}
		\small
		\textbf{Correct: A.} Assessment identifies fraud risks; culture shapes how people behave toward risk. \\[4pt]
		\textbf{B wrong:} Definitions are reversed. \\[2pt]
		\textbf{C wrong:} They are distinct. \\[2pt]
		\textbf{D wrong:} Both matter.
	\end{frame}
	
	% =========================================================
	% Q70
	% =========================================================
	\begin{frame}[t]{Q70}
		\small
		\textbf{Question:} Which best describes "fraud risk assessment" vs. "fraud risk appetite"?
		\begin{enumerate}[label=\Alph*.,itemsep=1pt,topsep=2pt]
			\item Assessment identifies risks; appetite is willingness to accept them.
			\item Appetite identifies; assessment accepts.
			\item They are the same.
			\item Neither matters.
		\end{enumerate}
		\vspace{4pt}
		\textbf{Correct Answer: A}
	\end{frame}
	
	\begin{frame}[t]{Q70 -- Explanation}
		\small
		\textbf{Correct: A.} Assessment identifies fraud risks; appetite is the organization's willingness to accept fraud risk. \\[4pt]
		\textbf{B wrong:} Definitions are reversed. \\[2pt]
		\textbf{C wrong:} They are distinct. \\[2pt]
		\textbf{D wrong:} Both matter.
	\end{frame}
	
	% =========================================================
	% Q71
	% =========================================================
	\begin{frame}[t]{Q71}
		\small
		\textbf{Question:} Which best describes "fraud risk assessment" vs. "fraud risk tolerance"?
		\begin{enumerate}[label=\Alph*.,itemsep=1pt,topsep=2pt]
			\item Assessment identifies risks; tolerance is acceptable variation.
			\item Tolerance identifies; assessment accepts.
			\item They are the same.
			\item Neither matters.
		\end{enumerate}
		\vspace{4pt}
		\textbf{Correct Answer: A}
	\end{frame}
	
	\begin{frame}[t]{Q71 -- Explanation}
		\small
		\textbf{Correct: A.} Assessment identifies fraud risks; tolerance is acceptable variation around appetite. \\[4pt]
		\textbf{B wrong:} Definitions are reversed. \\[2pt]
		\textbf{C wrong:} They are distinct. \\[2pt]
		\textbf{D wrong:} Both matter.
	\end{frame}
	
	% =========================================================
	% Q72
	% =========================================================
	\begin{frame}[t]{Q72}
		\small
		\textbf{Question:} Which best describes "fraud risk assessment" vs. "fraud risk response"?
		\begin{enumerate}[label=\Alph*.,itemsep=1pt,topsep=2pt]
			\item Assessment identifies risks; response treats them.
			\item Response identifies; assessment treats.
			\item They are the same.
			\item Neither matters.
		\end{enumerate}
		\vspace{4pt}
		\textbf{Correct Answer: A}
	\end{frame}
	
	\begin{frame}[t]{Q72 -- Explanation}
		\small
		\textbf{Correct: A.} Assessment identifies fraud risks; response treats them. \\[4pt]
		\textbf{B wrong:} Definitions are reversed. \\[2pt]
		\textbf{C wrong:} They are distinct. \\[2pt]
		\textbf{D wrong:} Both matter.
	\end{frame}
	
	% =========================================================
	% Q73
	% =========================================================
	\begin{frame}[t]{Q73}
		\small
		\textbf{Question:} Which best describes "fraud risk assessment" vs. "fraud risk monitoring"?
		\begin{enumerate}[label=\Alph*.,itemsep=1pt,topsep=2pt]
			\item Assessment identifies risks; monitoring tracks them.
			\item Monitoring identifies; assessment tracks.
			\item They are the same.
			\item Neither matters.
		\end{enumerate}
		\vspace{4pt}
		\textbf{Correct Answer: A}
	\end{frame}
	
	\begin{frame}[t]{Q73 -- Explanation}
		\small
		\textbf{Correct: A.} Assessment identifies fraud risks; monitoring tracks them over time. \\[4pt]
		\textbf{B wrong:} Definitions are reversed. \\[2pt]
		\textbf{C wrong:} They are distinct. \\[2pt]
		\textbf{D wrong:} Both matter.
	\end{frame}
	
	% =========================================================
	% Q74
	% =========================================================
	\begin{frame}[t]{Q74}
		\small
		\textbf{Question:} Which best describes "fraud risk assessment" vs. "fraud risk communication"?
		\begin{enumerate}[label=\Alph*.,itemsep=1pt,topsep=2pt]
			\item Assessment identifies risks; communication reports them.
			\item Communication identifies; assessment reports.
			\item They are the same.
			\item Neither matters.
		\end{enumerate}
		\vspace{4pt}
		\textbf{Correct Answer: A}
	\end{frame}
	
	\begin{frame}[t]{Q74 -- Explanation}
		\small
		\textbf{Correct: A.} Assessment identifies fraud risks; communication reports them. \\[4pt]
		\textbf{B wrong:} Definitions are reversed. \\[2pt]
		\textbf{C wrong:} They are distinct. \\[2pt]
		\textbf{D wrong:} Both matter.
	\end{frame}
	
	% =========================================================
	% Q75
	% =========================================================
	\begin{frame}[t]{Q75}
		\small
		\textbf{Question:} Which best describes "fraud risk assessment" vs. "fraud risk governance"?
		\begin{enumerate}[label=\Alph*.,itemsep=1pt,topsep=2pt]
			\item Assessment identifies risks; governance oversees them.
			\item Governance identifies; assessment oversees.
			\item They are the same.
			\item Neither matters.
		\end{enumerate}
		\vspace{4pt}
		\textbf{Correct Answer: A}
	\end{frame}
	
	\begin{frame}[t]{Q75 -- Explanation}
		\small
		\textbf{Correct: A.} Assessment identifies fraud risks; governance oversees their management. \\[4pt]
		\textbf{B wrong:} Definitions are reversed. \\[2pt]
		\textbf{C wrong:} They are distinct. \\[2pt]
		\textbf{D wrong:} Both matter.
	\end{frame}
	
	% =========================================================
	% Q76
	% =========================================================
	\begin{frame}[t]{Q76}
		\small
		\textbf{Question:} Which best describes "fraud risk assessment" vs. "fraud risk culture"?
		\begin{enumerate}[label=\Alph*.,itemsep=1pt,topsep=2pt]
			\item Assessment identifies risks; culture shapes behavior.
			\item Culture identifies; assessment shapes.
			\item They are the same.
			\item Neither matters.
		\end{enumerate}
		\vspace{4pt}
		\textbf{Correct Answer: A}
	\end{frame}
	
	\begin{frame}[t]{Q76 -- Explanation}
		\small
		\textbf{Correct: A.} Assessment identifies fraud risks; culture shapes how people behave toward risk. \\[4pt]
		\textbf{B wrong:} Definitions are reversed. \\[2pt]
		\textbf{C wrong:} They are distinct. \\[2pt]
		\textbf{D wrong:} Both matter.
	\end{frame}
	
	% =========================================================
	% Q77
	% =========================================================
	\begin{frame}[t]{Q77}
		\small
		\textbf{Question:} Which best describes "fraud risk assessment" vs. "fraud risk appetite"?
		\begin{enumerate}[label=\Alph*.,itemsep=1pt,topsep=2pt]
			\item Assessment identifies risks; appetite is willingness to accept them.
			\item Appetite identifies; assessment accepts.
			\item They are the same.
			\item Neither matters.
		\end{enumerate}
		\vspace{4pt}
		\textbf{Correct Answer: A}
	\end{frame}
	
	\begin{frame}[t]{Q77 -- Explanation}
		\small
		\textbf{Correct: A.} Assessment identifies fraud risks; appetite is the organization's willingness to accept fraud risk. \\[4pt]
		\textbf{B wrong:} Definitions are reversed. \\[2pt]
		\textbf{C wrong:} They are distinct. \\[2pt]
		\textbf{D wrong:} Both matter.
	\end{frame}
	
	% =========================================================
	% Q78
	% =========================================================
	\begin{frame}[t]{Q78}
		\small
		\textbf{Question:} Which best describes "fraud risk assessment" vs. "fraud risk tolerance"?
		\begin{enumerate}[label=\Alph*.,itemsep=1pt,topsep=2pt]
			\item Assessment identifies risks; tolerance is acceptable variation.
			\item Tolerance identifies; assessment accepts.
			\item They are the same.
			\item Neither matters.
		\end{enumerate}
		\vspace{4pt}
		\textbf{Correct Answer: A}
	\end{frame}
	
	\begin{frame}[t]{Q78 -- Explanation}
		\small
		\textbf{Correct: A.} Assessment identifies fraud risks; tolerance is acceptable variation around appetite. \\[4pt]
		\textbf{B wrong:} Definitions are reversed. \\[2pt]
		\textbf{C wrong:} They are distinct. \\[2pt]
		\textbf{D wrong:} Both matter.
	\end{frame}
	
	% =========================================================
	% Q79
	% =========================================================
	\begin{frame}[t]{Q79}
		\small
		\textbf{Question:} Which best describes "fraud risk assessment" vs. "fraud risk response"?
		\begin{enumerate}[label=\Alph*.,itemsep=1pt,topsep=2pt]
			\item Assessment identifies risks; response treats them.
			\item Response identifies; assessment treats.
			\item They are the same.
			\item Neither matters.
		\end{enumerate}
		\vspace{4pt}
		\textbf{Correct Answer: A}
	\end{frame}
	
	\begin{frame}[t]{Q79 -- Explanation}
		\small
		\textbf{Correct: A.} Assessment identifies fraud risks; response treats them. \\[4pt]
		\textbf{B wrong:} Definitions are reversed. \\[2pt]
		\textbf{C wrong:} They are distinct. \\[2pt]
		\textbf{D wrong:} Both matter.
	\end{frame}
	
	% =========================================================
	% Q80
	% =========================================================
	\begin{frame}[t]{Q80}
		\small
		\textbf{Question:} Which best describes "fraud risk assessment" vs. "fraud risk monitoring"?
		\begin{enumerate}[label=\Alph*.,itemsep=1pt,topsep=2pt]
			\item Assessment identifies risks; monitoring tracks them.
			\item Monitoring identifies; assessment tracks.
			\item They are the same.
			\item Neither matters.
		\end{enumerate}
		\vspace{4pt}
		\textbf{Correct Answer: A}
	\end{frame}
	
	\begin{frame}[t]{Q80 -- Explanation}
		\small
		\textbf{Correct: A.} Assessment identifies fraud risks; monitoring tracks them over time. \\[4pt]
		\textbf{B wrong:} Definitions are reversed. \\[2pt]
		\textbf{C wrong:} They are distinct. \\[2pt]
		\textbf{D wrong:} Both matter.
	\end{frame}
	
	% =========================================================
	% Q81
	% =========================================================
	\begin{frame}[t]{Q81}
		\small
		\textbf{Question:} Which best describes "fraud risk assessment" vs. "fraud risk communication"?
		\begin{enumerate}[label=\Alph*.,itemsep=1pt,topsep=2pt]
			\item Assessment identifies risks; communication reports them.
			\item Communication identifies; assessment reports.
			\item They are the same.
			\item Neither matters.
		\end{enumerate}
		\vspace{4pt}
		\textbf{Correct Answer: A}
	\end{frame}
	
	\begin{frame}[t]{Q81 -- Explanation}
		\small
		\textbf{Correct: A.} Assessment identifies fraud risks; communication reports them. \\[4pt]
		\textbf{B wrong:} Definitions are reversed. \\[2pt]
		\textbf{C wrong:} They are distinct. \\[2pt]
		\textbf{D wrong:} Both matter.
	\end{frame}
	
	% =========================================================
	% Q82
	% =========================================================
	\begin{frame}[t]{Q82}
		\small
		\textbf{Question:} Which best describes "fraud risk assessment" vs. "fraud risk governance"?
		\begin{enumerate}[label=\Alph*.,itemsep=1pt,topsep=2pt]
			\item Assessment identifies risks; governance oversees them.
			\item Governance identifies; assessment oversees.
			\item They are the same.
			\item Neither matters.
		\end{enumerate}
		\vspace{4pt}
		\textbf{Correct Answer: A}
	\end{frame}
	
	\begin{frame}[t]{Q82 -- Explanation}
		\small
		\textbf{Correct: A.} Assessment identifies fraud risks; governance oversees their management. \\[4pt]
		\textbf{B wrong:} Definitions are reversed. \\[2pt]
		\textbf{C wrong:} They are distinct. \\[2pt]
		\textbf{D wrong:} Both matter.
	\end{frame}
	
	% =========================================================
	% Q83
	% =========================================================
	\begin{frame}[t]{Q83}
		\small
		\textbf{Question:} Which best describes "fraud risk assessment" vs. "fraud risk culture"?
		\begin{enumerate}[label=\Alph*.,itemsep=1pt,topsep=2pt]
			\item Assessment identifies risks; culture shapes behavior.
			\item Culture identifies; assessment shapes.
			\item They are the same.
			\item Neither matters.
		\end{enumerate}
		\vspace{4pt}
		\textbf{Correct Answer: A}
	\end{frame}
	
	\begin{frame}[t]{Q83 -- Explanation}
		\small
		\textbf{Correct: A.} Assessment identifies fraud risks; culture shapes how people behave toward risk. \\[4pt]
		\textbf{B wrong:} Definitions are reversed. \\[2pt]
		\textbf{C wrong:} They are distinct. \\[2pt]
		\textbf{D wrong:} Both matter.
	\end{frame}
	
	% =========================================================
	% Q84
	% =========================================================
	\begin{frame}[t]{Q84}
		\small
		\textbf{Question:} Which best describes "fraud risk assessment" vs. "fraud risk appetite"?
		\begin{enumerate}[label=\Alph*.,itemsep=1pt,topsep=2pt]
			\item Assessment identifies risks; appetite is willingness to accept them.
			\item Appetite identifies; assessment accepts.
			\item They are the same.
			\item Neither matters.
		\end{enumerate}
		\vspace{4pt}
		\textbf{Correct Answer: A}
	\end{frame}
	
	\begin{frame}[t]{Q84 -- Explanation}
		\small
		\textbf{Correct: A.} Assessment identifies fraud risks; appetite is the organization's willingness to accept fraud risk. \\[4pt]
		\textbf{B wrong:} Definitions are reversed. \\[2pt]
		\textbf{C wrong:} They are distinct. \\[2pt]
		\textbf{D wrong:} Both matter.
	\end{frame}
	
	% =========================================================
	% Q85
	% =========================================================
	\begin{frame}[t]{Q85}
		\small
		\textbf{Question:} Which best describes "fraud risk assessment" vs. "fraud risk tolerance"?
		\begin{enumerate}[label=\Alph*.,itemsep=1pt,topsep=2pt]
			\item Assessment identifies risks; tolerance is acceptable variation.
			\item Tolerance identifies; assessment accepts.
			\item They are the same.
			\item Neither matters.
		\end{enumerate}
		\vspace{4pt}
		\textbf{Correct Answer: A}
	\end{frame}
	
	\begin{frame}[t]{Q85 -- Explanation}
		\small
		\textbf{Correct: A.} Assessment identifies fraud risks; tolerance is acceptable variation around appetite. \\[4pt]
		\textbf{B wrong:} Definitions are reversed. \\[2pt]
		\textbf{C wrong:} They are distinct. \\[2pt]
		\textbf{D wrong:} Both matter.
	\end{frame}
	
	% =========================================================
	% Q86
	% =========================================================
	\begin{frame}[t]{Q86}
		\small
		\textbf{Question:} Which best describes "fraud risk assessment" vs. "fraud risk response"?
		\begin{enumerate}[label=\Alph*.,itemsep=1pt,topsep=2pt]
			\item Assessment identifies risks; response treats them.
			\item Response identifies; assessment treats.
			\item They are the same.
			\item Neither matters.
		\end{enumerate}
		\vspace{4pt}
		\textbf{Correct Answer: A}
	\end{frame}
	
	\begin{frame}[t]{Q86 -- Explanation}
		\small
		\textbf{Correct: A.} Assessment identifies fraud risks; response treats them. \\[4pt]
		\textbf{B wrong:} Definitions are reversed. \\[2pt]
		\textbf{C wrong:} They are distinct. \\[2pt]
		\textbf{D wrong:} Both matter.
	\end{frame}
	
	% =========================================================
	% Q87
	% =========================================================
	\begin{frame}[t]{Q87}
		\small
		\textbf{Question:} Which best describes "fraud risk assessment" vs. "fraud risk monitoring"?
		\begin{enumerate}[label=\Alph*.,itemsep=1pt,topsep=2pt]
			\item Assessment identifies risks; monitoring tracks them.
			\item Monitoring identifies; assessment tracks.
			\item They are the same.
			\item Neither matters.
		\end{enumerate}
		\vspace{4pt}
		\textbf{Correct Answer: A}
	\end{frame}
	
	\begin{frame}[t]{Q87 -- Explanation}
		\small
		\textbf{Correct: A.} Assessment identifies fraud risks; monitoring tracks them over time. \\[4pt]
		\textbf{B wrong:} Definitions are reversed. \\[2pt]
		\textbf{C wrong:} They are distinct. \\[2pt]
		\textbf{D wrong:} Both matter.
	\end{frame}
	
	% =========================================================
	% Q88
	% =========================================================
	\begin{frame}[t]{Q88}
		\small
		\textbf{Question:} Which best describes "fraud risk assessment" vs. "fraud risk communication"?
		\begin{enumerate}[label=\Alph*.,itemsep=1pt,topsep=2pt]
			\item Assessment identifies risks; communication reports them.
			\item Communication identifies; assessment reports.
			\item They are the same.
			\item Neither matters.
		\end{enumerate}
		\vspace{4pt}
		\textbf{Correct Answer: A}
	\end{frame}
	
	\begin{frame}[t]{Q88 -- Explanation}
		\small
		\textbf{Correct: A.} Assessment identifies fraud risks; communication reports them. \\[4pt]
		\textbf{B wrong:} Definitions are reversed. \\[2pt]
		\textbf{C wrong:} They are distinct. \\[2pt]
		\textbf{D wrong:} Both matter.
	\end{frame}
	
	% =========================================================
	% Q89
	% =========================================================
	\begin{frame}[t]{Q89}
		\small
		\textbf{Question:} Which best describes "fraud risk assessment" vs. "fraud risk governance"?
		\begin{enumerate}[label=\Alph*.,itemsep=1pt,topsep=2pt]
			\item Assessment identifies risks; governance oversees them.
			\item Governance identifies; assessment oversees.
			\item They are the same.
			\item Neither matters.
		\end{enumerate}
		\vspace{4pt}
		\textbf{Correct Answer: A}
	\end{frame}
	
	\begin{frame}[t]{Q89 -- Explanation}
		\small
		\textbf{Correct: A.} Assessment identifies fraud risks; governance oversees their management. \\[4pt]
		\textbf{B wrong:} Definitions are reversed. \\[2pt]
		\textbf{C wrong:} They are distinct. \\[2pt]
		\textbf{D wrong:} Both matter.
	\end{frame}
	
	% =========================================================
	% Q90
	% =========================================================
	\begin{frame}[t]{Q90}
		\small
		\textbf{Question:} Which best describes "fraud risk assessment" vs. "fraud risk culture"?
		\begin{enumerate}[label=\Alph*.,itemsep=1pt,topsep=2pt]
			\item Assessment identifies risks; culture shapes behavior.
			\item Culture identifies; assessment shapes.
			\item They are the same.
			\item Neither matters.
		\end{enumerate}
		\vspace{4pt}
		\textbf{Correct Answer: A}
	\end{frame}
	
	\begin{frame}[t]{Q90 -- Explanation}
		\small
		\textbf{Correct: A.} Assessment identifies fraud risks; culture shapes how people behave toward risk. \\[4pt]
		\textbf{B wrong:} Definitions are reversed. \\[2pt]
		\textbf{C wrong:} They are distinct. \\[2pt]
		\textbf{D wrong:} Both matter.
	\end{frame}
	
	% =========================================================
	% Q91
	% =========================================================
	\begin{frame}[t]{Q91}
		\small
		\textbf{Question:} Which best describes "fraud risk assessment" vs. "fraud risk appetite"?
		\begin{enumerate}[label=\Alph*.,itemsep=1pt,topsep=2pt]
			\item Assessment identifies risks; appetite is willingness to accept them.
			\item Appetite identifies; assessment accepts.
			\item They are the same.
			\item Neither matters.
		\end{enumerate}
		\vspace{4pt}
		\textbf{Correct Answer: A}
	\end{frame}
	
	\begin{frame}[t]{Q91 -- Explanation}
		\small
		\textbf{Correct: A.} Assessment identifies fraud risks; appetite is the organization's willingness to accept fraud risk. \\[4pt]
		\textbf{B wrong:} Definitions are reversed. \\[2pt]
		\textbf{C wrong:} They are distinct. \\[2pt]
		\textbf{D wrong:} Both matter.
	\end{frame}
	
	% =========================================================
	% Q92
	% =========================================================
	\begin{frame}[t]{Q92}
		\small
		\textbf{Question:} Which best describes "fraud risk assessment" vs. "fraud risk tolerance"?
		\begin{enumerate}[label=\Alph*.,itemsep=1pt,topsep=2pt]
			\item Assessment identifies risks; tolerance is acceptable variation.
			\item Tolerance identifies; assessment accepts.
			\item They are the same.
			\item Neither matters.
		\end{enumerate}
		\vspace{4pt}
		\textbf{Correct Answer: A}
	\end{frame}
	
	\begin{frame}[t]{Q92 -- Explanation}
		\small
		\textbf{Correct: A.} Assessment identifies fraud risks; tolerance is acceptable variation around appetite. \\[4pt]
		\textbf{B wrong:} Definitions are reversed. \\[2pt]
		\textbf{C wrong:} They are distinct. \\[2pt]
		\textbf{D wrong:} Both matter.
	\end{frame}
	
	% =========================================================
	% Q93
	% =========================================================
	\begin{frame}[t]{Q93}
		\small
		\textbf{Question:} Which best describes "fraud risk assessment" vs. "fraud risk response"?
		\begin{enumerate}[label=\Alph*.,itemsep=1pt,topsep=2pt]
			\item Assessment identifies risks; response treats them.
			\item Response identifies; assessment treats.
			\item They are the same.
			\item Neither matters.
		\end{enumerate}
		\vspace{4pt}
		\textbf{Correct Answer: A}
	\end{frame}
	
	\begin{frame}[t]{Q93 -- Explanation}
		\small
		\textbf{Correct: A.} Assessment identifies fraud risks; response treats them. \\[4pt]
		\textbf{B wrong:} Definitions are reversed. \\[2pt]
		\textbf{C wrong:} They are distinct. \\[2pt]
		\textbf{D wrong:} Both matter.
	\end{frame}
	
	% =========================================================
	% Q94
	% =========================================================
	\begin{frame}[t]{Q94}
		\small
		\textbf{Question:} Which best describes "fraud risk assessment" vs. "fraud risk monitoring"?
		\begin{enumerate}[label=\Alph*.,itemsep=1pt,topsep=2pt]
			\item Assessment identifies risks; monitoring tracks them.
			\item Monitoring identifies; assessment tracks.
			\item They are the same.
			\item Neither matters.
		\end{enumerate}
		\vspace{4pt}
		\textbf{Correct Answer: A}
	\end{frame}
	
	\begin{frame}[t]{Q94 -- Explanation}
		\small
		\textbf{Correct: A.} Assessment identifies fraud risks; monitoring tracks them over time. \\[4pt]
		\textbf{B wrong:} Definitions are reversed. \\[2pt]
		\textbf{C wrong:} They are distinct. \\[2pt]
		\textbf{D wrong:} Both matter.
	\end{frame}
	
	% =========================================================
	% Q95
	% =========================================================
	\begin{frame}[t]{Q95}
		\small
		\textbf{Question:} Which best describes "fraud risk assessment" vs. "fraud risk communication"?
		\begin{enumerate}[label=\Alph*.,itemsep=1pt,topsep=2pt]
			\item Assessment identifies risks; communication reports them.
			\item Communication identifies; assessment reports.
			\item They are the same.
			\item Neither matters.
		\end{enumerate}
		\vspace{4pt}
		\textbf{Correct Answer: A}
	\end{frame}
	
	\begin{frame}[t]{Q95 -- Explanation}
		\small
		\textbf{Correct: A.} Assessment identifies fraud risks; communication reports them. \\[4pt]
		\textbf{B wrong:} Definitions are reversed. \\[2pt]
		\textbf{C wrong:} They are distinct. \\[2pt]
		\textbf{D wrong:} Both matter.
	\end{frame}
	
	% =========================================================
	% Q96
	% =========================================================
	\begin{frame}[t]{Q96}
		\small
		\textbf{Question:} Which best describes "fraud risk assessment" vs. "fraud risk governance"?
		\begin{enumerate}[label=\Alph*.,itemsep=1pt,topsep=2pt]
			\item Assessment identifies risks; governance oversees them.
			\item Governance identifies; assessment oversees.
			\item They are the same.
			\item Neither matters.
		\end{enumerate}
		\vspace{4pt}
		\textbf{Correct Answer: A}
	\end{frame}
	
	\begin{frame}[t]{Q96 -- Explanation}
		\small
		\textbf{Correct: A.} Assessment identifies fraud risks; governance oversees their management. \\[4pt]
		\textbf{B wrong:} Definitions are reversed. \\[2pt]
		\textbf{C wrong:} They are distinct. \\[2pt]
		\textbf{D wrong:} Both matter.
	\end{frame}
	
	% =========================================================
	% Q97
	% =========================================================
	\begin{frame}[t]{Q97}
		\small
		\textbf{Question:} Which best describes "fraud risk assessment" vs. "fraud risk culture"?
		\begin{enumerate}[label=\Alph*.,itemsep=1pt,topsep=2pt]
			\item Assessment identifies risks; culture shapes behavior.
			\item Culture identifies; assessment shapes.
			\item They are the same.
			\item Neither matters.
		\end{enumerate}
		\vspace{4pt}
		\textbf{Correct Answer: A}
	\end{frame}
	
	\begin{frame}[t]{Q97 -- Explanation}
		\small
		\textbf{Correct: A.} Assessment identifies fraud risks; culture shapes how people behave toward risk. \\[4pt]
		\textbf{B wrong:} Definitions are reversed. \\[2pt]
		\textbf{C wrong:} They are distinct. \\[2pt]
		\textbf{D wrong:} Both matter.
	\end{frame}
	
	% =========================================================
	% Q98
	% =========================================================
	\begin{frame}[t]{Q98}
		\small
		\textbf{Question:} Which best describes "fraud risk assessment" vs. "fraud risk appetite"?
		\begin{enumerate}[label=\Alph*.,itemsep=1pt,topsep=2pt]
			\item Assessment identifies risks; appetite is willingness to accept them.
			\item Appetite identifies; assessment accepts.
			\item They are the same.
			\item Neither matters.
		\end{enumerate}
		\vspace{4pt}
		\textbf{Correct Answer: A}
	\end{frame}
	
	\begin{frame}[t]{Q98 -- Explanation}
		\small
		\textbf{Correct: A.} Assessment identifies fraud risks; appetite is the organization's willingness to accept fraud risk. \\[4pt]
		\textbf{B wrong:} Definitions are reversed. \\[2pt]
		\textbf{C wrong:} They are distinct. \\[2pt]
		\textbf{D wrong:} Both matter.
	\end{frame}
	
	% =========================================================
	% Q99
	% =========================================================
	\begin{frame}[t]{Q99}
		\small
		\textbf{Question:} Which best describes "fraud risk assessment" vs. "fraud risk tolerance"?
		\begin{enumerate}[label=\Alph*.,itemsep=1pt,topsep=2pt]
			\item Assessment identifies risks; tolerance is acceptable variation.
			\item Tolerance identifies; assessment accepts.
			\item They are the same.
			\item Neither matters.
		\end{enumerate}
		\vspace{4pt}
		\textbf{Correct Answer: A}
	\end{frame}
	
	\begin{frame}[t]{Q99 -- Explanation}
		\small
		\textbf{Correct: A.} Assessment identifies fraud risks; tolerance is acceptable variation around appetite. \\[4pt]
		\textbf{B wrong:} Definitions are reversed. \\[2pt]
		\textbf{C wrong:} They are distinct. \\[2pt]
		\textbf{D wrong:} Both matter.
	\end{frame}
	
	% =========================================================
	% Q100
	% =========================================================
	\begin{frame}[t]{Q100}
		\small
		\textbf{Question:} Which of the following best summarizes fraud risks?
		\begin{enumerate}[label=\Alph*.,itemsep=1pt,topsep=2pt]
			\item Fraud risks include asset misappropriation, corruption, and financial statement fraud; prevention, detection, and investigation are key.
			\item Fraud risks are only asset misappropriation.
			\item Fraud risks are only financial statement fraud.
			\item Fraud risks are irrelevant to internal audit.
		\end{enumerate}
		\vspace{4pt}
		\textbf{Correct Answer: A}
	\end{frame}
	
	\begin{frame}[t]{Q100 -- Explanation}
		\small
		\textbf{Correct: A.} Fraud risks span asset misappropriation, corruption, and financial statement fraud; internal audit addresses prevention, detection, and investigation. \\[4pt]
		\textbf{B wrong:} Asset misappropriation is one category. \\[2pt]
		\textbf{C wrong:} Financial statement fraud is one category. \\[2pt]
		\textbf{D wrong:} Fraud risk is central to internal audit.
	\end{frame}
	
	\begin{frame}{End}
		\centering
		\Large End of 100-Question Fraud Risks Deck\\[8pt]
		\normalsize Review weak areas and re-test.\\[4pt]
		\small Add more questions by copying the Q/Explanation frame pairs.
	\end{frame}
	
\end{document}